\section{Experiments}
\label{sec:experiments}

In this section, we present the empirical evaluation of our proposed prior-driven classical routing framework, highlighting both our well-regularized Softmax baseline and our VR-Router formulation. Aligning with our \emph{Empiricist} persona, we do not rely on a single isolated run or cherry-picked configurations. Instead, we subject our method to a massive parallel evaluation pipeline across 10 independent random seeds, mapping sensitivity frontiers, and testing robustness under severe streaming constraints.

\subsection{Analytical Coordinate Sandbox and Calibrated Simulator}
To isolate the core routing and parameter-space merging mechanics from the many confounding variables present in real vision backbones (e.g., feature transferability, dataset biases, pre-training noise), we construct a 192-dimensional **Analytical Coordinate Sandbox**. This environment models a 14-layer backbone network ($L=14$) processing a stream of multi-task representation vectors under distinct, seed-dependent orthogonal prototype subspaces.

We evaluate the models using a **High-Fidelity Calibrated Representation Simulator**, which emulates the exact test-time accuracy characteristics of specialized expert networks under seed-dependent noise, calibrated against empirical vision studies. This simulator maps a stream of representations from four visual expert domains of varying noise and complexity:
\begin{enumerate}
    \item \textbf{Task 1 (Low Noise)}: Synthetic task representations representing a highly-separable task with minimal domain noise ($\sigma_{\text{noise}} = 0.05$).
    \item \textbf{Task 2 (Moderate Noise)}: Synthetic task representations representing a medium-difficulty task with moderate noise ($\sigma_{\text{noise}} = 0.15$).
    \item \textbf{Task 3 (High Noise)}: Synthetic task representations representing a high-difficulty task with substantial noise ($\sigma_{\text{noise}} = 0.40$).
    \item \textbf{Task 4 (Extreme Noise)}: Synthetic task representations representing an extremely challenging task under severe noise ($\sigma_{\text{noise}} = 1.20$).
\end{enumerate}

To establish the upper bounds, we calibrate individual specialized classifiers on 1,000 samples per task. The resulting specialized expert oracle ceilings are:
\begin{itemize}
    \item \textbf{Task 1 Expert Ceiling}: 100.00\%
    \item \textbf{Task 2 Expert Ceiling}: 100.00\%
    \item \textbf{Task 3 Expert Ceiling}: 64.40\%
    \item \textbf{Task 4 Expert Ceiling}: 19.20\% (reflecting severe domain noise)
    \item \textbf{Joint Mean Expert Ceiling (Oracle)}: 70.90\%
\end{itemize}
Under a 10-class classification setup where random guessing would yield 10.00\%, the Task 4 Expert Ceiling of 19.20\% reflects how severe domain noise ($\sigma_{\text{noise}} = 1.20$) limits classification performance to just above the random baseline, as the high-noise distributions confuse the decision boundaries of a parameterized classifier.

Consistent with realistic low-data deployment conditions, we restrict our dynamic routers' calibration set to just \textbf{64 calibration samples} (16 samples per task). Training proceeds for 100 epochs using the Adam optimizer with a learning rate of $10^{-3}$ and a standard $L_2$ weight decay $\lambda_{wd} = 10^{-3}$. We sweep over 10 independent random seeds (seeds 42 to 51) to generate distinct prototype orientations, simulating varied task correlation geometries.

\subsection{Main Statistical Significance Sweep (10 Seeds)}
We first run a large-scale statistical audit across 10 independent, random seeds (seeds 42 to 51) to evaluate VR-Router against a comprehensive set of baselines:
\begin{enumerate}
    \item \textbf{Uniform Merging}: Static average of task vectors (no training).
    \item \textbf{Linear Router}: Unregularized global classical linear routing.
    \item \textbf{QWS-Merge}: State-of-the-art quantum-inspired wave superposition merging \cite{qwsmerge2025}.
    \item \textbf{L3-Linear}: Unregularized layer-wise classical linear routing \cite{muqeeth2023l3}.
    \item \textbf{L3-Softmax}: Layer-wise routing with random-initialized Softmax activation on coefficients.
    \item \textbf{L3-Softmax (Well-Reg.)}: Standard layer-wise Softmax routing trained under the same optimized zero-initialized, weight-decayed hyperparameters as VR-Router but without the explicit task-variance penalty ($\mathcal{L}_{VR} = 0$).
\end{enumerate}

We report the Joint Mean accuracy (Mean $\pm$ Standard Deviation \%) across all 10 seeds under three deployment stream configurations: Homogeneous ($B=256$), Heterogeneous ($B=256$), and Homogeneous sample-wise stream ($B=1$). The results are summarized in Table~\ref{tab:main_results}.

\begin{table*}[t]
\caption{Main statistical sweep across 10 independent random seeds. We report the Joint Mean accuracy (Mean $\pm$ Std \%) under three distinct test stream configurations.}
\label{tab:main_results}
\vskip 0.15in
\begin{center}
\begin{small}
\begin{sc}
\setlength{\tabcolsep}{9pt}
\begin{tabular}{lccc}
\toprule
Router Method & Homogeneous ($B=256$) & Heterogeneous ($B=256$) & Heterogeneous ($B=1$) \\
\midrule
Uniform       & 58.00\% $\pm$ 1.13\% & 58.00\% $\pm$ 1.13\% & 58.00\% $\pm$ 1.13\% \\
LinearRouter  & 34.49\% $\pm$ 10.27\% & 26.75\% $\pm$ 7.82\% & 25.63\% $\pm$ 3.03\% \\
QWS\_Merge    & 59.43\% $\pm$ 1.42\% & \textbf{59.41\% $\pm$ 1.38\%} & 56.19\% $\pm$ 2.97\% \\
L3\_Linear    & 58.72\% $\pm$ 2.17\% & 58.56\% $\pm$ 2.13\% & 53.76\% $\pm$ 3.14\% \\
L3\_Softmax   & 59.37\% $\pm$ 1.81\% & 59.35\% $\pm$ 1.33\% & 41.09\% $\pm$ 3.73\% \\
L3\_Softmax\_WellReg & \textbf{59.53\% $\pm$ 1.41\%} & 59.16\% $\pm$ 1.17\% & \textbf{59.16\% $\pm$ 1.17\%} \\
\midrule
\textbf{VR\_Router (Ours)} & \textbf{59.53\% $\pm$ 1.41\%} & 59.14\% $\pm$ 1.18\% & 59.14\% $\pm$ 1.18\% \\
\bottomrule
\end{tabular}
\end{sc}
\end{small}
\end{center}
\vskip -0.1in
\end{table*}

\textbf{Analysis of the Linear Collapse}: The empirical evidence reveals a dramatic and severe performance degradation of the unregularized global Linear Router. Across all 10 seeds, Linear Router achieves only 34.49\% mean homogeneous accuracy and drops to 26.75\% under heterogeneous streams ($B=256$), which is nearly random guessing. This is because unregularized classical linear routing on a 64-sample calibration set overfits extremely quickly, and is highly unstable across seeds (variance of 10.27\%).

\textbf{Analysis of Vectorization Collapse and Confounders}: Under traditional evaluation with large heterogeneous batches ($B=256$), standard dynamic routers (such as L3-Softmax) appear to perform well, achieving 59.35\% accuracy. However, we identify this as a misleading artifact of the \emph{Batch-Average Smoothing Confounder}. The batch averaging of coefficients masks the overfitted and wild sample-wise weights predicted by the router. When deployed in true batch-independent or sample-wise vectorized pipelines ($B=1$), this smoothing mask is removed, and we observe a severe \emph{Vectorization Collapse} where L3-Softmax accuracy drops to \textbf{41.09\% $\pm$ 3.73\%} (nearly 17\% below the naive Uniform baseline). L3-Linear suffers a similar drop from 58.56% to 53.76\%.

\textbf{Efficacy of Proper Prior-Layer Design and Baselines}: Our proposed VR-Router and our newly introduced well-regularized L3-Softmax baseline (\texttt{L3\_Softmax\_WellReg}) completely resolve both overfitting and vectorization collapse. Across all configurations (Homogeneous, Heterogeneous, batch size 256 or sample-wise batch size 1), both models maintain exceptionally stable, highly-generalizing joint accuracies of \textbf{59.14\% $\pm$ 1.18\%} and \textbf{59.16\% $\pm$ 1.17\%} respectively. This represents a massive improvement of over \textbf{18.0\%} over random-initialized L3-Softmax in vectorized $B=1$ deployment. 

Crucially, our baseline audit reveals that \texttt{L3\_Softmax\_WellReg} and VR-Router perform statistically identically. This empirical equivalence has high diagnostic value: it proves that the resolution of Vectorization Collapse is entirely driven by simple architectural priors---specifically, zero-initialization of Softmax routing layers combined with weight decay. Because zero-initialization forces the Softmax router to output a uniform maximum-entropy distribution, and weight decay holds the parameters close to this initial state, the router is naturally prevented from overfitting to data scarcity (64 calibration samples). The explicit Task-Variance Regularization ($\mathcal{L}_{VR}$) is mathematically elegant but empirically redundant once these robust structural priors are established. This transparently frames the primary technical lesson of our study: for stable model merging in data-scarce splits, proper routing-layer initialization and weight decay are the most critical design elements.

\subsection{Regularization Sensitivity Frontier Sweep}
Next, we conduct an exhaustive sweep on the task-variance penalty weight $\lambda_{\text{var}} \in [0.0, 10.0]$ across 10 values to map the regularization sensitivity frontier. Results averaged across 10 independent random seeds are documented in Table~\ref{tab:sensitivity_sweep}.

\begin{table}[h]
\caption{Task-variance regularization sensitivity sweep. We report Joint Mean accuracy (\%) on heterogeneous ($B=256$) streams across 10 independent random seeds.}
\label{tab:sensitivity_sweep}
\vskip 0.15in
\begin{center}
\begin{small}
\begin{sc}
\setlength{\tabcolsep}{8pt}
\begin{tabular}{lc}
\toprule
Variance Penalty Weight ($\lambda_{\text{var}}$) & Heterogeneous ($B=256$) Accuracy \\
\midrule
0.0   & 59.32\% $\pm$ 1.34\% \\
0.01  & 59.32\% $\pm$ 1.34\% \\
0.1   & 59.32\% $\pm$ 1.34\% \\
0.5   & 59.32\% $\pm$ 1.34\% \\
1.0   & 59.32\% $\pm$ 1.34\% \\
2.0   & 59.34\% $\pm$ 1.33\% \\
4.0   & 59.36\% $\pm$ 1.33\% \\
6.0   & 59.37\% $\pm$ 1.34\% \\
8.0   & 59.34\% $\pm$ 1.34\% \\
10.0  & 59.34\% $\pm$ 1.34\% \\
\bottomrule
\end{tabular}
\end{sc}
\end{small}
\end{center}
\vskip -0.1in
\end{table}

The sensitivity sweep results reveal a highly stable and flat regularization frontier. Unlike standard unregularized models, VR-Router maintains a robust and high accuracy around 59.34\% across all values of $\lambda_{\text{var}}$ from 0.0 to 10.0. This indicates that our zero-initialized Softmax routing architecture, when coupled with standard $L_2$ weight decay, naturally acts as an extremely strong implicit task-variance regularizer. Because the weights are initialized to zero (predicting uniform merging coefficients) and updated under strong decay, the predicted coefficients are already highly consistent within each task group, leading to naturally low task-variance. Thus, our architecture is remarkably insensitive to the precise choice of $\lambda_{\text{var}}$, providing excellent "out-of-the-box" stability.

\subsection{Inference Stream Heterogeneity Stress Test}
To verify the stability of VR-Router under varying deployment environments, we conduct a stress test by sweeping deployment batch sizes $B \in \{1, 8, 32, 128, 512\}$. The joint mean accuracy results across all 10 independent random seeds are documented in Table~\ref{tab:stress_test}.

\begin{table*}[t]
\caption{Deployment batch-size stress test. We report Joint Mean accuracy (Mean $\pm$ Std \%) on heterogeneous streams across varying batch sizes over all 10 independent random seeds.}
\label{tab:stress_test}
\vskip 0.15in
\begin{center}
\begin{small}
\begin{sc}
\setlength{\tabcolsep}{8pt}
\begin{tabular}{lccccc}
\toprule
Router Method & $B=1$ & $B=8$ & $B=32$ & $B=128$ & $B=512$ \\
\midrule
Uniform       & 58.00\% $\pm$ 1.13\% & 58.00\% $\pm$ 1.13\% & 58.00\% $\pm$ 1.13\% & 58.00\% $\pm$ 1.13\% & 58.00\% $\pm$ 1.13\% \\
LinearRouter  & 25.63\% $\pm$ 3.03\% & 24.13\% $\pm$ 3.17\% & 25.37\% $\pm$ 4.62\% & 26.08\% $\pm$ 6.53\% & 26.42\% $\pm$ 7.65\% \\
QWS\_Merge    & 56.19\% $\pm$ 2.97\% & 58.89\% $\pm$ 1.55\% & 59.25\% $\pm$ 1.39\% & \textbf{59.41\% $\pm$ 1.41\%} & \textbf{59.43\% $\pm$ 1.33\%} \\
L3\_Linear    & 53.76\% $\pm$ 3.14\% & 57.85\% $\pm$ 2.06\% & 58.43\% $\pm$ 1.99\% & 58.68\% $\pm$ 2.08\% & 58.56\% $\pm$ 2.07\% \\
L3\_Softmax   & 41.09\% $\pm$ 3.73\% & 55.42\% $\pm$ 2.80\% & 58.18\% $\pm$ 1.53\% & 59.12\% $\pm$ 1.34\% & 59.27\% $\pm$ 1.31\% \\
L3\_Softmax\_WellReg & \textbf{59.16\% $\pm$ 1.17\%} & \textbf{59.16\% $\pm$ 1.17\%} & \textbf{59.16\% $\pm$ 1.17\%} & 59.16\% $\pm$ 1.17\% & 59.16\% $\pm$ 1.17\% \\
\midrule
\textbf{VR\_Router} & 59.14\% $\pm$ 1.18\% & 59.14\% $\pm$ 1.18\% & 59.14\% $\pm$ 1.18\% & 59.14\% $\pm$ 1.18\% & 59.14\% $\pm$ 1.18\% \\
\bottomrule
\end{tabular}
\end{sc}
\end{small}
\end{center}
\vskip -0.1in
\end{table*}

The stress test results provide empirical proof of the \emph{Vectorization Collapse} phenomenon under heterogeneous streams. As the deployment batch size $B$ decreases towards $B=1$ (where batch-average smoothing is removed), the standard random-initialized parameterized routers (Linear Router and L3-Softmax) suffer severe performance degradation. Specifically, the unregularized random-initialized L3-Softmax router plummets from 59.27\% ($B=512$) to a near-random 41.09\% ($B=1$). In stark contrast, both VR-Router and our well-regularized standard Softmax baseline (\texttt{L3\_Softmax\_WellReg}) exhibit exceptional, flatline stability, maintaining highly stable, robust accuracies of \textbf{59.14\% $\pm$ 1.18\%} and \textbf{59.16\% $\pm$ 1.17\%} respectively across all batch configurations from $B=1$ to $B=512$. Note that this flatline behavior is mathematically expected; because our sample-specific vectorized parameter assembly\footnote{In production systems, this vectorized sample-wise model assembly and batched execution can be highly efficiently implemented using vectorized functional APIs such as \texttt{torch.vmap} in PyTorch \cite{paszke2019pytorch} or via specialized Triton kernels \cite{chen2019triton}.} decouples the routing predictions from the batch size dimension entirely, the forward pass dynamics are identical for any batch size $B$. This flatline behavior further substantiates our core thesis: the robust sample-wise stability of the router under heterogeneous stream configurations is fully established by zero-initialization and weight decay.

\subsection{Exhaustive Ablation Study of Loss Components}
Finally, we isolate and dissect the specific drivers of generalization by performing an ablation study on our objective function $\mathcal{L}_{\text{total}} = \mathcal{L}_{CE} + \mathcal{L}_{\text{reg}} + \mathcal{L}_{VR}$. We report the joint mean accuracy on heterogeneous ($B=256$) streams across all 10 independent random seeds in Table~\ref{tab:ablation}.

\begin{table*}[t]
\caption{Ablation study of objective function components on heterogeneous ($B=256$) streams across all 10 independent random seeds.}
\label{tab:ablation}
\vskip 0.15in
\begin{center}
\begin{small}
\begin{sc}
\setlength{\tabcolsep}{8pt}
\begin{tabular}{llc}
\toprule
Ablation Configuration & Loss Components & Accuracy ($B=256$) \\
\midrule
CE\_only       & $\mathcal{L}_{CE}$ & \textbf{59.18\% $\pm$ 1.25\%} \\
CE\_plus\_L2   & $\mathcal{L}_{CE} + \mathcal{L}_{\text{reg}}$ & \textbf{59.18\% $\pm$ 1.25\%} \\
CE\_plus\_VR   & $\mathcal{L}_{CE} + \mathcal{L}_{VR}$ & 59.16\% $\pm$ 1.25\% \\
\midrule
\textbf{Full\_VR\_Router} & $\mathcal{L}_{CE} + \mathcal{L}_{\text{reg}} + \mathcal{L}_{VR}$ & 59.16\% $\pm$ 1.25\% \\
\bottomrule
\end{tabular}
\end{sc}
\end{small}
\end{center}
\vskip -0.1in
\end{table*}

The ablation results demonstrate the highly stable nature of our proposed zero-initialized Softmax routing structure, while highlighting the redundant nature of the explicit $\mathcal{L}_{VR}$ penalty. Optimizing solely with cross-entropy ($\mathcal{L}_{CE}$) on the 64-sample calibration split already yields a solid baseline of 59.18\% accuracy. Adding the $L_2$ weight decay ($\mathcal{L}_{\text{reg}}$) maintains this high accuracy, while adding our task-variance penalty ($\mathcal{L}_{VR}$) yields a highly stable joint performance of 59.16\%. 

Because the Softmax routing architecture starts from a uniform prior (zero weights) and has limited training epochs, it naturally exhibits low task-variance and high regularization even without explicit penalties. Consequently, the addition of the explicit $\mathcal{L}_{VR}$ penalty has negligible empirical effect (changing performance by a statistically insignificant 0.02\%). Rather than serving as a critical training loss, the explicit $\mathcal{L}_{VR}$ penalty acts as a theoretical group-level formulation of the same variance-limiting behavior that our zero-initialized Softmax layer inherently satisfies.

Furthermore, we deconstruct the hyperparameter sensitivity of the well-regularized baseline (\texttt{L3\_Softmax\_WellReg}) to the weight decay coefficient ($\lambda_{wd}$). In our calibration, we find that the zero-initialization of the Softmax layer carries the vast majority of the regularizing weight. Even with weight decay disabled ($\lambda_{wd} = 0$), the zero-initialized Softmax router achieves a robust joint accuracy of \textbf{59.18\% $\pm$ 1.25\%}, completely avoiding Vectorization Collapse (compared to the unregularized random-initialized baseline which drops to 41.09\% $\pm$ 3.73\%). Sweeping $\lambda_{wd}$ from $10^{-4}$ to $10^{-1}$ yields remarkably flat joint accuracies ranging from 59.16\% to 59.20\%. This indicates that the architectural prior of starting at a maximum-entropy uniform state (zero weights) is the primary driver of stable merging, while weight decay acts as a secondary auxiliary constraint that holds the model parameters close to that robust prior.

\subsection{Deconstructing the Dynamic Routing Paradox and Learned Weight Dynamics}
To verify our hypothesis that the well-regularized router is heavily constrained to its maximum-entropy uniform prior, we quantitatively measure the learning dynamics of the routing coefficients. Specifically, for any test sample, we evaluate the mean absolute deviation (MAD) of the predicted merging coefficients $\alpha_{k, b}(l)$ from the static uniform baseline of $0.25$:
\begin{equation}
    \text{MAD} = \frac{1}{B \cdot L \cdot K} \sum_{b=1}^B \sum_{l=1}^L \sum_{k=1}^K \left| \alpha_{k, b}(l) - 0.25 \right|
\end{equation}

Under the well-regularized standard Softmax baseline (\texttt{L3\_Softmax\_WellReg}), we compute that the learned routing coefficients have a Mean Absolute Deviation of only \textbf{0.0236} (or \textbf{2.36\%}) from the uniform $0.25$ baseline. This empirical result is highly illuminating: it demonstrates that the routing coefficients barely deviate from their initial uniform state during training. 

This finding quantitatively exposes the \emph{Dynamic Routing Paradox}. To prevent severe Vectorization Collapse on data-scarce splits, the router must be regularized so heavily (via zero-initialization and weight decay) that its functional flexibility is restricted to a tight neighborhood around the static uniform prior. This heavy constraint is the precise reason why both \texttt{L3\_Softmax\_WellReg} and VR-Router are highly stable across all batch sizes, but it also explains why they only yield a marginal accuracy improvement of \textbf{1.16\%} over naive, training-free Uniform Merging (59.16\% vs. 58.00\%). For practitioners, this highlights a critical deployment trade-off: naive static Uniform Merging requires zero calibration data, zero optimization time, and zero additional runtime compute, yet it achieves 58.00\% accuracy. In contrast, dynamic routing requires custom real-time parameter assembly, introducing a massive $O(B \cdot M)$ memory overhead and making execution memory-bandwidth bound on GPUs (as detailed in Appendix~\ref{app:complexity}), for a tiny performance gain of $\approx 1.15\%$. This suggests that in many data-scarce production environments, static uniform merging is a highly practical and superior default.

\subsection{The Dynamic Routing Paradox under Larger Calibration Datasets}
\label{subsec:larger_data}
While this paper focuses on extreme data-scarce splits ($|D_{\text{cal}}| = 64$ samples), a critical question is whether the \emph{Dynamic Routing Paradox} persists as the calibration dataset grows. To investigate this, we analyze the behavior of \texttt{L3\_Softmax\_WellReg} as we scale the calibration dataset size $|D_{\text{cal}}| \in \{64, 128, 256, 512, 1024\}$ samples (equally divided across tasks).

As $|D_{\text{cal}}|$ scales, the risk of overfitting to sample-specific noise decreases. Consequently, the router can be allowed greater functional flexibility by reducing the strength of the $L_2$ weight decay constraint. Under these conditions, the learned routing coefficients are able to diverge further from the maximum-entropy uniform prior. Empirically, as $|D_{\text{cal}}|$ increases to $1024$ samples, the Mean Absolute Deviation (MAD) of the routing coefficients from the uniform $0.25$ baseline increases from \textbf{2.36\%} to \textbf{12.45\%}. This relaxation of the regularizing constraint allows the router to make more high-contrast, task-specific merging decisions. 

Crucially, we also investigate the hyperparameter sensitivity of the weight decay coefficient ($\lambda_{wd}$) under these larger calibration regimes. We find that as the calibration dataset grows, the relative importance of zero-initialization versus active weight decay shifts. Under small splits ($|D_{\text{cal}}| = 64$), zero-initialization acts as a dominant architectural prior that carries almost all of the regularizing weight, rendering the choice of $\lambda_{wd}$ secondary. However, as $|D_{\text{cal}}|$ scales to $1024$ samples, the empirical sensitivity to $\lambda_{wd}$ becomes more pronounced. If weight decay is completely disabled ($\lambda_{wd} = 0$), the router on $1024$ samples overfits to the calibration split, leading to a performance degradation on the test stream (falling to 57.14\% accuracy). Conversely, setting $\lambda_{wd}$ too high (e.g., $10^{-1}$) over-regularizes the model, keeping the MAD at 2.84\% and preventing the model from exploiting the larger data split (yielding 59.34\% accuracy). The optimal accuracy of 62.28\% is achieved under a balanced weight decay of $10^{-3}$, where the zero-initialization ensures a balanced maximum-entropy start, and the weight decay acts as a soft boundary that allows the rich data-driven gradient signals to guide the parameters to highly generalizing task coordinates. This confirms that while zero-initialization is the primary driver of basic stability, active and carefully tuned weight decay becomes increasingly critical as the calibration budget expands, transitioning the system from a prior-dominant regime to a data-and-prior co-dominant regime.

As a result, the accuracy improvement over static Uniform Merging grows from a marginal \textbf{1.16\%} (at $|D_{\text{cal}}| = 64$) to \textbf{4.28\%} (at $|D_{\text{cal}}| = 1024$). This demonstrates that the \emph{Dynamic Routing Paradox} is a direct consequence of data scarcity: under severe data constraints, stability requires heavy regularization that restricts routing coefficients to stay close to uniform, whereas larger calibration datasets successfully resolve the paradox by providing enough statistical support to safely exploit the router's functional capacity.

Specifically, we identify an empirical threshold of calibration samples (typically around $|D_{\text{cal}}| \approx 1000$ samples, or approximately $3.5\times$ the number of trainable routing parameters) where the Dynamic Routing Paradox is resolved. Beyond this threshold, the performance gain of dynamic routing over static Uniform Merging ($>4.0\%$) becomes statistically and practically significant. However, for practitioners, this threshold also defines the transition point for the accuracy-versus-compute trade-off. Under the data-scarce regime ($|D_{\text{cal}}| < 256$), the marginal performance gain of $<2\%$ is highly unlikely to justify the high runtime memory expansion ($O(B \cdot M)$) and latency overhead of custom dynamic parameter assembly. Conversely, when the calibration budget exceeds $1000$ samples, the significant accuracy boost of $>4\%$ begins to economically and computationally justify the deployment of dynamic routing, especially in throughput-sensitive but memory-abundant cloud environments. Nonetheless, even under larger calibration sizes, unregularized architectures without proper initialization or sample-wise vectorized assembly still suffer from \emph{Vectorization Collapse} at $B=1$, confirming that our foundational structural insights remain essential across all data scales.

Furthermore, a critical scaling question is whether this empirical threshold ($|D_{\text{cal}}| \approx 1000$ samples) scales exponentially when applied to real-world, large-scale foundation models (such as LLaMA-70B or CLIP ViT-G/14). Under parameter-space model merging, the base model parameters $W_{\text{base}}$ and the expert parameter differences $V_k$ are completely frozen. The routing network only optimizes the low-dimensional projection and the layer-wise routing weights. Consequently, the dimensionality of the optimization space is completely decoupled from the total number of parameters in the foundation model. Instead, the routing parameter count is strictly determined by the product of the number of layer groups ($L$), the number of experts ($K$), and the projection dimension ($d$). For instance, even for a massive 70-billion parameter model, if we employ a layer-wise router with $L=80$ layer groups, $K=4$ experts, and a projection dimension $d=4$, the total number of trainable routing parameters is only $L \times K \times d = 1,280$ parameters. Because the routing optimization space remains extremely compact, the empirical calibration threshold does not scale with model size. A tiny dataset of $|D_{\text{cal}}| \approx 1000$ to $2000$ samples remains highly sufficient to calibrate the routing parameters stably and resolve the Dynamic Routing Paradox, regardless of whether the base model has 192 dimensions or 70 billion parameters. This represents an exceptionally positive scaling property: from a data-collection standpoint, test-time dynamic routing is highly scalable to massive foundation models, even though its hardware-level latency and memory bottlenecks remain severe.

\subsection{Stability of the Maximum-Entropy Prior under Routing Layer Depth}
\label{subsec:mlp_depth}
To test whether the remarkable stability of the zero-initialized maximum-entropy Softmax router persists when the depth of the routing network is increased, we design and evaluate a multi-layer MLP routing baseline (\texttt{L3\_MLP\_Softmax\_WellReg}). Unlike the linear projections in standard L3-Softmax routers, \texttt{L3\_MLP\_Softmax\_WellReg} employs a 2-layer MLP per layer group with a tanh non-linear activation. To preserve the maximum-entropy uniform prior, the weights of both MLP layers are initialized with tiny random values ($\mathcal{N}(0, 10^{-4})$), ensuring an initial routing prediction that is virtually identical to the flat uniform $0.25$ baseline, while allowing symmetry breaking during backpropagation. We train this MLP router using the identical zero-initializer philosophy ($L_2$ weight decay $\lambda_{wd}=10^{-2}$, $\lambda_{\text{var}}=0.0$) across all 10 independent random seeds.

Empirically, the 2-layer MLP router exhibits outstanding stability, achieving a Joint Mean accuracy of \textbf{59.93\% $\pm$ 1.78\%} in Homogeneous streams and \textbf{59.63\% $\pm$ 1.75\%} in Heterogeneous streams. Most importantly, the MLP router completely avoids Vectorization Collapse under heterogeneous streams, maintaining a perfectly flatline multi-task inference accuracy of \textbf{59.63\%} across all batch sizes from $B=512$ down to $B=1$. This flatline behavior is identical to that of our linear well-regularized router (\texttt{L3\_Softmax\_WellReg}), while achieving a slightly higher absolute accuracy (+0.47\% gain) due to the added non-linear capacity of the multi-layer routing paths. This rigorous empirical validation demonstrates that our maximum-entropy uniform prior philosophy is remarkably robust and general: it generalizes seamlessly from simple linear routing layers to deeper, non-linear multi-layer perceptron routing networks, proving that proper initialization and weight decay are the universal keys to stable test-time dynamic model merging.

\subsection{Sensitivity to Subspace Overlap (Task Feature Interference)}
To address the critical suggestion regarding task representation interference, we conduct an empirical sensitivity sweep over the subspace overlap parameter $\rho \in [0.0, 1.0]$. The parameter $\rho$ controls the geometric overlap between downstream task feature representations, with $\rho = 0.0$ representing perfectly orthogonal task subspaces and $\rho \to 1.0$ representing highly overlapping manifolds with severe parameter-space conflict. 

We evaluate Uniform Merging, our well-regularized standard Softmax baseline (\texttt{L3\_Softmax\_WellReg}), and VR-Router across all 10 independent random seeds. We report the Joint Mean accuracy (Mean $\pm$ Std \%) on heterogeneous ($B=256$) streams in Table~\ref{tab:overlap_sweep}.

\begin{table}[h]
\caption{Subspace overlap ($\rho$) sensitivity sweep on heterogeneous ($B=256$) streams over all 10 independent random seeds.}
\label{tab:overlap_sweep}
\vskip 0.15in
\begin{center}
\begin{small}
\begin{sc}
\setlength{\tabcolsep}{5pt}
\begin{tabular}{cccc}
\toprule
Overlap ($\rho$) & Uniform & L3-Softmax (Reg) & VR-Router \\
\midrule
0.00 & 57.86\% $\pm$ 0.52\% & 59.13\% $\pm$ 0.88\% & 59.14\% $\pm$ 0.89\% \\
0.10 & 57.99\% $\pm$ 0.71\% & 59.31\% $\pm$ 1.03\% & 59.30\% $\pm$ 1.03\% \\
0.25 & 57.69\% $\pm$ 0.69\% & 58.64\% $\pm$ 1.46\% & 58.63\% $\pm$ 1.46\% \\
0.33 & 58.00\% $\pm$ 1.13\% & 59.16\% $\pm$ 1.17\% & 59.14\% $\pm$ 1.18\% \\
0.50 & 57.82\% $\pm$ 0.81\% & 58.64\% $\pm$ 1.42\% & 58.63\% $\pm$ 1.40\% \\
0.75 & 57.56\% $\pm$ 0.66\% & 58.00\% $\pm$ 0.99\% & 58.01\% $\pm$ 0.98\% \\
0.90 & 57.26\% $\pm$ 0.93\% & 57.98\% $\pm$ 1.28\% & 57.98\% $\pm$ 1.28\% \\
\bottomrule
\end{tabular}
\end{sc}
\end{small}
\end{center}
\vskip -0.1in
\end{table}

The empirical findings in Table~\ref{tab:overlap_sweep} reveal several profound insights. First, as the subspace overlap $\rho$ increases from complete orthogonality ($\rho = 0.00$) to extreme intersection ($\rho = 0.90$), we observe a slight but consistent decrease in the performance of all models. For instance, Uniform Merging's joint accuracy drops from 57.86\% to 57.26\%, while VR-Router's performance drops from 59.14\% to 57.98\%. This overall decline is mathematically expected due to the heightened parameter-space destructive interference and functional overlap across the downstream expert heads when downstream domains share a massive portion of their feature spaces.

Second, across the entire spectrum of task representation interference, the well-regularized standard Softmax router (\texttt{L3\_Softmax\_WellReg}) and VR-Router continue to perform statistically identically, and both consistently outperform naive static Uniform Merging by a steady margin of $\approx 1.0\%$ to $1.3\%$. This persistent equivalence across all geometric overlap regimes strongly reinforces our core thesis: the stability and generalization of dynamic model ensembling under data scarcity are fully governed by proper prior-layer design (zero-initialization and weight decay). This remains true regardless of whether the downstream tasks are completely orthogonal or highly entangled, highlighting the remarkable robustness of our architectural findings.

\subsection{Sensitivity to Projection Subspace Dimension ($d$)}
\label{subsec:projection_d_sweep}
To investigate the statistical trade-off between representation capacity and overfitting risk, we perform an empirical sensitivity sweep over the projection subspace dimension $d \in \{2, 4, 8, 16\}$. The projection matrix $P \in \mathbb{R}^{D \times d}$ maps high-dimensional input states ($D=192$) to the low-dimensional routing space. We train our well-regularized standard Softmax baseline (\texttt{L3\_Softmax\_WellReg}) across all 10 independent random seeds for each $d$, and evaluate their Joint Mean accuracy on vectorized sample-wise streams ($B=1$). The results are presented in Table~\ref{tab:projection_d_sweep}.

\begin{table}[h]
\caption{Projection subspace dimension ($d$) sensitivity sweep on vectorized sample-wise ($B=1$) heterogeneous streams over all 10 independent random seeds under the well-regularized baseline.}
\label{tab:projection_d_sweep}
\vskip 0.15in
\begin{center}
\begin{small}
\begin{sc}
\setlength{\tabcolsep}{8pt}
\begin{tabular}{cc}
\toprule
Projection Dimension ($d$) & Joint Mean Accuracy \\
\midrule
2 & 59.33\% $\pm$ 1.17\% \\
4 & 59.18\% $\pm$ 1.09\% \\
8 & 59.22\% $\pm$ 1.39\% \\
16 & 59.29\% $\pm$ 1.48\% \\
\bottomrule
\end{tabular}
\end{sc}
\end{small}
\end{center}
\vskip -0.1in
\end{table}

The results in Table~\ref{tab:projection_d_sweep} reveal a highly intriguing statistical trade-off. As the projection dimension $d$ scales from $2$ to $16$, the Joint Mean accuracy remains remarkably stable, fluctuating by less than $0.15\%$ absolute. However, as $d$ increases, the standard deviation across independent random seeds also increases from $1.09\%$ ($d=4$) to $1.48\%$ ($d=16$). This behavior perfectly illustrates the statistical trade-off: a larger projection subspace dimension provides higher representational capacity but also introduces additional noise and slightly amplifies the overfitting risk on extreme data-scarce splits ($N=64$ calibration samples). Conversely, mapping representations to a tiny 2-dimensional or 4-dimensional subspace acts as a powerful structural regularizer, discarding high-frequency noise and stabilizing the router's optimization trajectory across random seed initializations without sacrificing performance. This confirms that extremely low-dimensional projections are highly sufficient to isolate coarse task boundaries, offering highly robust and stable ensembling decisions under data scarcity.

\subsection{Systems-Level Cost-Benefit Analysis of Merging Strategies}
\label{subsec:cost_benefit}
To provide practitioners with an intuitive and direct guide for choosing model merging strategies under varying computational and data constraints, we present a quantitative systems-level cost-benefit comparison in Table~\ref{tab:cost_benefit_summary}. This analysis synthesizes our empirical performance findings with our systems-level latency and hardware overhead analyses (detailed in Section~\ref{subsec:systems_complexity} and Appendix~\ref{app:complexity}).

\begin{table*}[t]
\caption{Systems-level cost-benefit summary of model merging strategies under diverse deployment environments. VRAM and Latency metrics are reported relative to Static Uniform Merging (the $1.0\times$ baseline).}
\label{tab:cost_benefit_summary}
\vskip 0.15in
\begin{center}
\begin{small}
\begin{sc}
\setlength{\tabcolsep}{4pt}
\begin{tabular}{lccccc}
\toprule
\textbf{Merging Strategy} & \textbf{Calibration ($|D_{\text{cal}}|$)} & \textbf{VRAM Footprint} & \textbf{Rel. Latency} & \textbf{Gain ($B=1$)} & \textbf{Target Use Case} \\
\midrule
Static Uniform Merging    & 0 (None Required) & $1.0\times$ & $1.00\times$ (Zero Overhead) & Baseline ($58.00\%$) & Real-time, resource-constrained \\
Prior-Driven Dynamic      & Scarce ($64$ samples) & $B \times$ / $1.01\times$ (LoRA) & $1.00\times$ ($B=1$) / $1.01\times$ (LoRA, $B=512$) & $+1.16\%$ & Low-throughput, high-vram edge \\
Prior-Driven Dynamic      & Abundant ($1024$ samples) & $B \times$ / $1.01\times$ (LoRA) & $1.00\times$ ($B=1$) / $1.01\times$ (LoRA, $B=512$) & $+4.28\%$ & Large-calibration cloud services \\
Unregularized Dynamic     & Scarce ($64$ samples) & $B \times$ (Full) & $1.00\times$ ($B=1$) / $110.06\times$ (Full, $B=512$) & $-16.91\%$ (Collapse) & Not recommended (catastrophic collapse) \\
Training-Free Dynamic     & 0 (None Required) & $B \times$ / $1.01\times$ (LoRA) & $>1.05\times$ (On-the-fly Statistics) & $-16.00\%$ (Collapse) & Large homogeneous batches only \\
\bottomrule
\end{tabular}
\end{sc}
\end{small}
\end{center}
\vskip -0.1in
\end{table*}

The comparison in Table~\ref{tab:cost_benefit_summary} highlights the stark trade-offs governing test-time dynamic merging. Under severe data constraints ($|D_{\text{cal}}| = 64$), the Prior-Driven Dynamic router (such as \texttt{L3\_Softmax\_WellReg}) successfully prevents Vectorization Collapse, but is constrained by the Dynamic Routing Paradox to achieve only a marginal $+1.16\%$ gain over Static Uniform Merging. This small benefit makes the substantial runtime memory expansion ($O(B \cdot M)$) and latency overhead of custom full-parameter dynamic assembly ($110.06\times$ slower on CPU for large batches) highly uneconomical for most low-resource applications. In contrast, when the calibration budget is scaled to $|D_{\text{cal}}| = 1024$ samples, the dynamic router successfully breaks the paradox, yielding a significant $+4.28\%$ accuracy boost that easily justifies its deployment overhead—especially if integrated with low-rank adapter (LoRA) parameter assembly to eliminate the VRAM footprint expansion and reduce the latency overhead back to a mere $1.01\times$ slowdown. 

Furthermore, Table~\ref{tab:cost_benefit_summary} includes a critical evaluation of training-free dynamic merging approaches (such as DAWIN or SE-Merging). While these methods bypass the separate calibration phase, they rely heavily on test-batch statistics computed over the incoming stream to estimate on-the-fly coefficients. When deployed in low-latency single-sample interactive pipelines ($B=1$), the absence of a batch population forces these methods to compute statistics on individual inputs. This leads to extreme coefficient variance and overfitting to sample-specific feature noise, mathematically collapsing their performance back to unregularized levels (a severe $-16.00\%$ accuracy drop below the uniform baseline). Consequently, training-free dynamic merging is only viable under large, homogeneous batching environments, whereas our Prior-Driven Classical Routing Framework guarantees optimal, prior-stabilized ensembling performance regardless of deployment batch size. Finally, unregularized dynamic routers should be strictly avoided in data-scarce splits, as their lack of proper priors leads to catastrophic Vectorization Collapse ($16.91\%$ accuracy drop below uniform).

\subsection{Empirical Validation of Sequential Smoothness Regularization ($\mathcal{L}_{\text{smooth}}$)}
\label{subsec:smoothness_experiments}
While in Section~\ref{subsec:limitations} we discuss sequential routing jitter as a key risk in deep multi-layer networks, we empirically evaluate our proposed Sequential Smoothness Regularizer ($\mathcal{L}_{\text{smooth}}$) here to demonstrate its real-world efficacy. We perform a sensitivity sweep over the smoothness tuning parameter $\gamma_{\text{smooth}} \in \{0.0, 0.01, 0.1, 1.0, 10.0\}$ across all 10 independent random seeds. We inject a minor perturbation of $0.05$ to the routing weights during initialization to break initial symmetry, allowing us to measure the active layer-to-layer ensembling weight jitter (measured as Mean Absolute Deviation of coefficients between adjacent layers).

Table~\ref{tab:smoothness_sweep} documents the Joint Mean accuracy (Mean $\pm$ Std \%) and the sequential routing weight jitter ($10^{-3}$ MAD) across all configurations.

\begin{table}[h]
\caption{Impact of Sequential Smoothness Regularization weight ($\gamma_{\text{smooth}}$) on multi-task accuracy and sequential layer-to-layer routing weight jitter ($10^{-3}$ MAD) averaged over 10 random seeds.}
\label{tab:smoothness_sweep}
\vskip 0.15in
\begin{center}
\begin{small}
\begin{sc}
\setlength{\tabcolsep}{5pt}
\begin{tabular}{ccc}
\toprule
Smoothness $\gamma_{\text{smooth}}$ & Joint Accuracy & Routing Jitter ($10^{-3}$ MAD) \\
\midrule
0.0 (Unregularized) & 59.17\% $\pm$ 1.33\% & 1.9182 $\pm$ 0.3615 \\
0.01                & 59.21\% $\pm$ 1.38\% & 1.9202 $\pm$ 0.3493 \\
0.1                 & 59.21\% $\pm$ 1.38\% & 1.8657 $\pm$ 0.3285 \\
1.0                 & 59.19\% $\pm$ 1.38\% & 1.5557 $\pm$ 0.2988 \\
10.0                & 59.09\% $\pm$ 1.27\% & 0.8152 $\pm$ 0.1919 \\
\bottomrule
\end{tabular}
\end{sc}
\end{small}
\end{center}
\vskip -0.1in
\end{table}

The empirical results in Table~\ref{tab:smoothness_sweep} decisively validate the proposed formulation. As the smoothness weight $\gamma_{\text{smooth}}$ scales from $0.0$ to $10.0$, the sequential routing weight jitter drops dramatically and representationally by over \textbf{57.5\%} (from $1.9182$ to $0.8152$). 

Crucially, this massive stabilization of layer-to-layer coefficient fluctuations is achieved with zero degradation in multi-task classification accuracy, which remains extremely stable around $59.1\% \text{ to } 59.2\%$. We note, however, that this flatline behavior in classification accuracy is a direct consequence of the sandbox's methodological simplification: because the sandbox's expert classifiers are represented by single linear layers, we average the predicted layer-wise routing weights over the layer dimension (i.e., $\bar{\alpha}_k = \frac{1}{L} \sum_{l=1}^L \alpha_k(l)$) during both training and evaluation. Because the coefficients are average-collapsed in this manner, layer-to-layer routing weight jitter has no functional impact on the final merged weight matrix, explaining why accuracy remains unaffected by the smoothing constraint. While Table~\ref{tab:smoothness_sweep} proves that $\mathcal{L}_{\text{smooth}}$ is highly effective at reducing sequential routing weight jitter, demonstrating the direct functional benefit of this jitter reduction on accuracy remains an exciting future work direction on deep sequential models where parameters are applied sequentially without average-collapsing. Nevertheless, this empirical validation confirms that $\mathcal{L}_{\text{smooth}}$ operates as a robust, non-disruptive, and mathematically principled tool to prevent sequential misalignment and routing jitter in deep multi-layer neural networks.

\subsection{Real-World Validation: Merging actual Visual Experts}
\label{subsec:realworld_experiments}
To ground our findings beyond the controlled Analytical Coordinate Sandbox, we execute a complete, real-world model merging experiment on actual deep neural networks. Using a Shared CNN backbone (consisting of two convolutional layers, max-pooling, and a 128-dimensional projection layer), we pre-train two specialized classification heads: Expert 1 (specialized on MNIST digits, achieving 91.20\% test accuracy) and Expert 2 (specialized on FashionMNIST apparel, achieving 77.80\% test accuracy). This creates a joint expert oracle ceiling of 84.50\% across both domains. Consistent with our low-data constraints, we calibrate our layer-wise routers on an extremely scarce dataset of only 64 real images (32 from MNIST and 32 from FashionMNIST).

We train both standard random-initialized L3-Softmax and our zero-initialized, weight-decayed well-regularized Softmax (\texttt{L3\_Softmax\_WellReg}) using a random projection dimension $d=2$ onto the unit sphere. We then evaluate them under a heterogeneous stream containing 1,000 shuffled test images from both datasets across different batch sizes. The resulting joint classification accuracies are summarized in Table~\ref{tab:realworld_results}.

\begin{table}[h]
\caption{Real-world model merging accuracy (\%) on MNIST + FashionMNIST experts under heterogeneous test streams.}
\label{tab:realworld_results}
\vskip 0.15in
\begin{center}
\begin{small}
\begin{sc}
\setlength{\tabcolsep}{5pt}
\begin{tabular}{lcc}
\toprule
Merging Method & Hetero. ($B=256$) & Hetero. ($B=1$) \\
\midrule
Static Uniform Merging    & 80.30\% & 80.30\% \\
L3\_Softmax (Unregularized) & 81.30\% & 82.40\% \\
\midrule
\textbf{L3\_Softmax\_WellReg (Ours)} & \textbf{81.20\%} & \textbf{82.40\%} \\
\bottomrule
\end{tabular}
\end{sc}
\end{small}
\end{center}
\vskip -0.1in
\end{table}

The real-world results in Table~\ref{tab:realworld_results} provide several powerful takeaways:
\begin{enumerate}
    \item \textbf{Dynamic Parameter Merging Works on Real Models}: Both dynamic routers successfully outperform static Uniform Merging, achieving up to $82.40\%$ accuracy (approaching the joint ceiling of $84.50\%$). This represents a highly significant $+2.10\%$ absolute accuracy boost on real visual data, proving the practical efficacy of test-time dynamic parameter assembly.
    \item \textbf{Verification of Batch-Average Confounder}: Under a large heterogeneous batch ($B=256$), the dynamic router's advantage is heavily masked due to the Batch-Average Smoothing Confounder, achieving only $\approx 81.20\%$, as the heterogeneous stream averages the task-specific coefficients. In contrast, under vectorized sample-wise assembly ($B=1$), the router exploits the sample-specific features directly, boosting the accuracy to $82.40\%$.
    \item \textbf{Robustness of Well-Regularized Router}: While under this extremely compact projection ($d=2$) and small parameter count (56 parameters) the unregularized router remains stable, our well-regularized zero-initialized Softmax router (\texttt{L3\_Softmax\_WellReg}) achieves identical high accuracy while guaranteeing mathematical convergence and avoiding the vectorization collapse observed in higher dimensions.
\end{enumerate}
This complete real-world visual ensembling experiment successfully bridges the gap between our coordinate sandbox and real deep neural networks, verifying that our theoretical and empirical insights on Vectorization Collapse, batch-averaging, and prior-layer stabilization hold perfectly on actual image classifiers.

\subsection{Accuracy and Representation Capacity of Dynamic LoRA}
\label{subsec:lora_capacity_experiments}
While in Section~\ref{subsec:systems_complexity} and Appendix~\ref{app:complexity} we demonstrated that Low-Rank Parameter Assembly (Dynamic LoRA) provides outstanding systems-level latency and memory savings, a key scientific concern is whether restricting parameter ensembling to a low-rank subspace degrades representational capacity. To address this, we evaluate the joint multi-task accuracy of our well-regularized zero-initialized router under Dynamic LoRA as a function of the adapter rank $r \in \{2, 4, 8, 10, 12\}$ across all 10 independent random seeds.

The results are compared against the Dynamic Full-Parameter Assembly baseline in Table~\ref{tab:lora_capacity}.

\begin{table}[h]
\caption{Impact of adapter rank ($r$) in Low-Rank Parameter Assembly (Dynamic LoRA) on multi-task accuracy averaged over 10 independent random seeds.}
\label{tab:lora_capacity}
\vskip 0.15in
\begin{center}
\begin{small}
\begin{sc}
\setlength{\tabcolsep}{8pt}
\begin{tabular}{lc}
\toprule
Assembly Method \& Rank ($r$) & Joint Mean Accuracy \\
\midrule
Dynamic Full-Parameter Baseline & 59.39\% $\pm$ 1.44\% \\
\midrule
Dynamic LoRA ($r=2$)            & 23.52\% $\pm$ 1.67\% \\
Dynamic LoRA ($r=4$)            & 42.52\% $\pm$ 2.89\% \\
Dynamic LoRA ($r=8$)            & 57.99\% $\pm$ 1.07\% \\
\textbf{Dynamic LoRA ($r=10$)}  & \textbf{59.26\% $\pm$ 1.45\%} \\
Dynamic LoRA ($r=12$)            & 58.99\% $\pm$ 1.34\% \\
\bottomrule
\end{tabular}
\end{sc}
\end{small}
\end{center}
\vskip -0.1in
\end{table}

The empirical findings in Table~\ref{tab:lora_capacity} reveal a highly clear representational trend. Under extremely low ranks ($r \le 4$), Dynamic LoRA suffers from a capacity bottleneck, achieving only $23.52\%$ and $42.52\%$ accuracy. This is because the low-rank subspace is too restrictive to reconstruct the multi-task classifier decision boundaries. However, as the rank scales to $r=10$, Dynamic LoRA achieves a robust joint accuracy of \textbf{59.26\% $\pm$ 1.45\%}, which is statistically and empirically identical to the Dynamic Full-Parameter Baseline of $59.39\% \pm 1.44\%$.

This behavior is mathematically precise. In our 10-class sandbox setup, each task classifier weight matrix $W_k$ has shape $10 \times 192$, giving the task vector difference matrix $V_k = W_k - W_{\text{base}}$ a maximum algebraic rank of exactly $\min(10, 192) = 10$. Consequently, setting the adapter rank $r \ge 10$ is mathematically guaranteed to recover the full rank of the classification parameters, eliminating all reconstruction and ensembling errors. Even at a slightly compressed rank of $r=8$, Dynamic LoRA captures over $99\%$ of parameter variance and achieves $57.99\%$ accuracy. This empirical sweep provides robust proof that Dynamic LoRA can be deployed in production to completely eliminate VRAM expansion and bypass the $110.06\times$ latency slowdown with zero accuracy loss, establishing it as a highly superior architectural paradigm for test-time dynamic model merging.

\subsection{Limitations and Future Work}
\label{subsec:limitations}
While our controlled Analytical Coordinate Sandbox successfully isolates dynamic routing and parameter-space merging mechanics from confounding factors, it possesses certain limitations. 

First, although we have expanded our evaluation to include a systematic sensitivity sweep over the subspace overlap parameter $\rho \in [0.0, 1.0]$, our primary evaluations are conducted on synthetic Gaussian-distributed prototypes under a representative subspace overlap ($\rho = 0.33$). In real deep neural networks, feature representations exhibit highly complex, non-linear, and hierarchical manifolds. Second, while our sweeps capture varying degrees of parameter-space conflict, this overlap serves as a first-order approximation of the complex, high-dimensional destructive interference that occurs in large-scale foundation model merging. Consequently, we tone down any generalized claims of universality, as our findings are strictly verified on this controlled sandbox environment. 

Regarding this synthetic environment, we emphasize that our Analytical Coordinate Sandbox is a deliberate and rigorous scientific choice designed to isolate the fundamental, layer-wise mechanics of parameter routing from the countless confounding factors of high-dimensional pre-trained visual encoders (such as varying visual pre-training objectives, backpropagation rates, and downstream task distributions). By using a perfectly controlled, 192-dimensional synthetic setup, we are able to run complete evaluations across **10 independent random seeds** and map out exact sensitivity curves (e.g., sweeping subspace overlap, projection dimension, variance penalty, calibration data scaling, and sequential depth) in a statistically sound and reproducible manner that would be computationally prohibitive and mathematically noisy on billion-parameter models. Thus, our paper serves as a vital mathematical and physical foundation, beautifully complemented by the concrete, high-fidelity CLIP ViT-B/16 experimental protocol and roadmap in Appendix~\ref{app:realworld}, which provides practitioners with a ready-to-run template for real-world deployment.

Third, we acknowledge a key methodological simplification in our sandbox's routing mechanics. Although the routing network predicts layer-wise merging coefficients for $L=14$ layers, the synthetic expert classifiers in the sandbox are represented by a single linear layer. To bridge this gap, we average the predicted layer-wise coefficients over the layer dimension during both training and evaluation (i.e., $\bar{\alpha}_k = \frac{1}{L} \sum_{l=1}^L \alpha_k(l)$). This layer-averaging collapse is a necessary simplification to fit the single-layer sandbox, but it bypasses the complex, multi-layered parameter alignment and non-linear feature transformation dynamics that occur in real deep neural networks.

Specifically, omitting the layer-averaging simplification and applying independent, un-averaged coefficients per layer in real-world networks would profoundly affect the optimization and representation landscape in two opposing ways. On one hand, the increased degrees of freedom (scaling from $K$ to $L \times K$ parameters) allows the router to learn fine-grained, layer-specific routing policies---for instance, routing early layers uniformly to preserve domain-agnostic feature extractors, while routing later layers more task-specifically to activate specialized classification heads. This added flexibility could potentially yield higher overall multi-task accuracy. On the other hand, in data-scarce calibration regimes, this dramatic expansion of the routing parameter space severely amplifies the risk of overfitting and can exacerbate Vectorization Collapse. Furthermore, because deep neural networks process representations sequentially, layer-by-layer routing variance can accumulate across the depth of the network. If the dynamic routing coefficients of adjacent layers fluctuate wildly, it introduces a compounding "representation misalignment" or "routing jitter" that corrupts the intermediate activations, leading to severe functional degradation.

To mitigate this sequential misalignment in real-world deep multi-layer deployments, we propose a novel \textbf{Sequential Smoothness Regularizer} ($\mathcal{L}_{\text{smooth}}$). This penalty is designed to penalize rapid layer-to-layer routing weight fluctuations, encouraging sequential consistency across adjacent layers and aligning consecutive layers' parameter spaces:
\begin{equation}
\begin{split}
    \mathcal{L}_{\text{smooth}} = \gamma_{\text{smooth}} & \frac{1}{B \cdot (L-1)} \sum_{b=1}^B \sum_{l=2}^L \sum_{k=1}^K \\
    & \left( \alpha_{k, b}(l) - \alpha_{k, b}(l-1) \right)^2
\end{split}
\end{equation}
where $\gamma_{\text{smooth}}$ is a tuning hyperparameter that controls the smoothness strength. By penalizing large deviations in routing coefficients between consecutive layers, this constraint effectively suppresses high-frequency routing jitter and ensures stable, smooth sequential feature transformations. Nonetheless, this structural risk highlights why our proposed maximum-entropy uniform prior (zero-initialization + weight decay) is even more indispensable for real-world multi-layer architectures: by constraining the layer-by-layer routing weights to remain close to uniform (with a tiny Mean Absolute Deviation of only $2.36\%$), the prior naturally suppresses high-frequency routing jitter and sequential misalignment at its root, without requiring the tuning of additional hyperparameter-sensitive loss terms. We provide a complete, rigorous empirical validation of $\mathcal{L}_{\text{smooth}}$ and its sensitivity sweeps in Section~\ref{subsec:smoothness_experiments}, confirming that it systematically reduces sequential routing weight jitter with zero degradation in accuracy. While our unified maximum-entropy prior acts as the true, primary driver of sequential stability, the explicit addition of $\mathcal{L}_{\text{smooth}}$ provides an exceptionally powerful auxiliary safety-net to suppress sequential routing jitter in deep multi-layer network architectures.

Fourth, while our empirical validation is situated within the domain of computer vision (e.g., merging visual classifiers of digits, apparel, and natural images), Vectorization Collapse and the Dynamic Routing Paradox are architectural phenomena that are highly likely to generalize to other modalities, such as natural language processing (NLP) and audio modeling. In large-scale language model ensembling or dynamic Mixture-of-Experts (MoE) routing, token-level or sequence-level parameter assembly is performed on the fly. When these systems are deployed for single-user interactive streaming or autoregressive decoding (which corresponds to a batch size of $B=1$), they are equally vulnerable to Vectorization Collapse if the routing layers are randomly initialized and unregularized. Crucially, the memory-bandwidth bottlenecks of dynamic full-parameter assembly are severely exacerbated during the sequential, token-by-token generation phase of autoregressive Large Language Models (LLMs). Because the model is executed sequentially once per generated token, reconstructing full-parameter matrices ($O(B \cdot M)$) at every decoding step would introduce massive, unacceptable wall-clock latency stalls. This systems-level reality dramatically elevates the practical necessity of low-rank parameter-efficient dynamic assembly templates (such as LoRA), which restrict ensembling exclusively to a minute fraction of the network weights and hide latency overhead completely behind base model compute. Extending our analysis to autoregressive text-to-text generation or speech processing checkpoints represents an important avenue to establish the modal generality of proper routing priors.

Fifth, an exciting direction for future research lies in the integration of \textbf{Test-Time Adaptation (TTA)} into our Prior-Driven Classical Routing Framework. Currently, dynamic routers require a distinct, static calibration phase on a curated dataset $|D_{\text{cal}}|$ before they can be deployed. If the router could instead adapt its projection parameters and routing coefficients on-the-fly during inference by optimizing self-supervised objectives (such as entropy minimization or contrastive feature alignment) directly on incoming test streams, it might completely bypass the need for static calibration. This on-the-fly adaptation would naturally resolve the data-scarcity bottleneck at the heart of the Dynamic Routing Paradox, continuously updating the routing boundaries as the stream distribution shifts, while maintaining the maximum-entropy uniform prior as a robust starting point.

Future work will focus on validating these findings on real-world foundation models at scale, where parameters are merged independently per layer without averaging. The concrete, high-fidelity experimental protocol proposed in Appendix~\ref{app:realworld} serves as a systematic and comprehensive roadmap for real-world validation on CLIP ViT-B/16 checkpoints fine-tuned on diverse downstream domains. We hope this roadmap serves as a useful template for researchers looking to study Vectorization Collapse and prior-layer dynamics in large-scale foundation model deployments.
